%% ME_QuestionPool_Additional_MC.tex
%% ─────────────────────────────────────────────────────────────────────────────
%% ADDITIONAL CONTROL QUESTIONS (Single Choice / type=mc)
%% Microeconomics, Jan Suchanek, Hochschule Bremen, Summer Semester 2026
%%
%% This file is a SUPPLEMENT to ME_QuestionPool_4_.tex. It contains only new
%% \begin{question}[type=mc]{...}...\end{question} blocks, grouped under the
%% same chapter headings as the master file, so each block can be copy-pasted
%% straight into the matching \section{...} of the main question pool.
%%
%% New material drawn from the script that was not yet covered by the existing
%% pool: Veblen goods, value chains, market delimitation, public/private goods
%% (Ch.1); Homo Economicus & simplifying assumptions (Ch.2); Law of One Price,
%% Hotelling's Law, economies of scale/DEVEX-CAPEX-OPEX, market forms grid,
%% PPF shape (Ch.3); prohibitive price & saturation quantity (Ch.4); natural
%% monopoly & antitrust regimes (Ch.6); the Fuel Crisis & Tax Rebate case study,
%% which currently has NO control questions at all (new Ch.7 section); patents,
%% risk vs. uncertainty, auctions (Ch.8); levels of the state, shell model,
%% government expenditure ratio (Ch.9); Stag Hunt & Chicken Game (Ch.10);
%% quantitative theory of money, functions of money, gains from trade (Ch.11);
%% non-tariff barriers (Ch.12).
%%
%% To compile this file standalone for proofreading, use the same preamble as
%% ME_QuestionPool_4_.tex (hsbquestions.sty, hsbcommon.sty must be present).
%% ─────────────────────────────────────────────────────────────────────────────

\documentclass[paper=a5]{scrartcl}
\usepackage{xcolor}
\usepackage{hsbcommon}
\usepackage[hidelinks]{hyperref}
\usepackage[utf8]{inputenc}
\usepackage{setspace}
\usepackage{pgfplots}
\pgfplotsset{compat=1.18}
\usepackage{hsbquestions}
\showanswerstrue                 %% ANSWER KEY view for proofreading — set to
                                  %% \showanswersfalse before distributing.

\title{Microeconomics}
\subtitle{Additional Control Questions (Single Choice) — Supplement}
\author{Claude (Anthropic) using the Script of Jan \noun{Suchanek}}

\begin{document}
\onehalfspacing
\maketitle
\clearpage

%% ═════════════════════════════════════════════════════════════════════════════
\section{Chapter 1 — Recap and Basics}
%% ═════════════════════════════════════════════════════════════════════════════

\begin{question}[type=mc]{Veblen Goods}
As the price of a luxury watch rises, demand for it \emph{also} rises, because
buyers value the watch partly as a status symbol. This good is best classified
as:
\begin{choices}
  \choice        A Giffen good.
  \choice        An inferior good.
  \correctchoice A Veblen good.
  \choice        A normal good with unit elasticity.
\end{choices}
\begin{solution}
  A \textbf{Veblen good} shows demand rising as price rises (like a Giffen
  good), but for a different reason: conspicuous consumption. The good itself
  is valued more highly \emph{because} it is expensive, e.g.\ luxury watches.
  A Giffen good rises in demand due to the income effect of a strongly
  inferior good, not because of status signalling.
\end{solution}
\end{question}

\begin{question}[type=mc]{Value Chain: Raw Materials to Merchandise}
Along the value chain from raw material to final product, which sequence is
correct?
\begin{choices}
  \choice        Merchandise $\to$ commodity $\to$ raw material.
  \correctchoice Raw material $\to$ commodity $\to$ merchandise.
  \choice        Commodity $\to$ raw material $\to$ merchandise.
  \choice        There is no fixed sequence; all three terms are synonyms.
\end{choices}
\begin{solution}
  Raw materials (e.g.\ crude oil, iron ore) are mined or grown, then
  \textbf{refined} into commodities (e.g.\ petrol, steel), which are then
  \textbf{manufactured} into merchandise. Raw materials are a subset of
  commodities, and commodities are a subset of merchandise — though this is
  not a strict legal distinction. Capital and labour are needed at every
  stage, and trade/logistics connect the stages.
\end{solution}
\end{question}

\begin{question}[type=mc]{Market Delimitation}
When economists say a market must be ``delimited'', they mean that one has to
clarify all of the following \emph{except}:
\begin{choices}
  \choice        Spatial delimitation (where the good is traded).
  \choice        Temporal delimitation (when it is traded).
  \choice        Substitutability on the buyer and seller side.
  \correctchoice The nationality of the company's shareholders.
\end{choices}
\begin{solution}
  Delimiting a market means specifying \emph{where} (spatial), \emph{when}
  (temporal), and \emph{by which characteristics} — in particular
  substitutability from the buyer's and the seller's perspective — the good
  or market is defined. Shareholder nationality is not part of market
  delimitation.
\end{solution}
\end{question}

\begin{question}[type=mc]{Public Goods and the Tragedy of the Commons}
A public good is defined as being:
\begin{choices}
  \choice        Excludable and rival.
  \correctchoice Non-excludable and non-rival (or, for common-pool resources,
                 non-excludable but rival).
  \choice        Always provided efficiently by private markets.
  \choice        Owned exclusively by the state.
\end{choices}
\begin{solution}
  \textbf{Public goods} (e.g.\ clean air, common knowledge) are
  \textbf{non-excludable}: nobody can be stopped from using them. Pure public
  goods are also \textbf{non-rival}. Common-pool resources (e.g.\ fish stocks,
  forests) are non-excludable but \emph{rival}, which is why they tend to be
  over-used — the tragedy of the commons. Private goods, by contrast, are
  excludable and more likely to be allocated efficiently through markets.
\end{solution}
\end{question}


%% ═════════════════════════════════════════════════════════════════════════════
\section{Chapter 2 — The Consumer}
%% ═════════════════════════════════════════════════════════════════════════════

\begin{question}[type=mc]{Homo Economicus and Diminishing Marginal Utility}
Which utility function best captures the idea that ``more is better than
less'', while also exhibiting \emph{diminishing} marginal utility?
\begin{choices}
  \choice        $U(x) = ax + b$
  \correctchoice $U(x) = \ln(x)$
  \choice        $U(x) = -x^2$
  \choice        $U(x) = \text{constant}$
\end{choices}
\begin{solution}
  A linear utility function $U(x) = ax+b$ satisfies $dU/dx > 0$ but has
  \emph{constant} marginal utility. A logarithmic function $U(x) = \ln(x)$
  is still increasing in $x$, but its marginal utility $dU/dx = 1/x$
  \textbf{falls} as $x$ rises — a more realistic depiction of diminishing
  returns to consumption.
\end{solution}
\end{question}

\begin{question}[type=mc]{Rational Behaviour in the Long Run}
The ``rational behaviour'' assumption of Homo Economicus is best understood
as meaning that:
\begin{choices}
  \choice        Every individual makes the utility-maximising decision at
                 every single moment.
  \correctchoice In the long run and under competitive pressure, actors who
                 do not act in their own best interest tend to be priced out
                 of the market.
  \choice        The state is subject to exactly the same competitive
                 pressures as private households and firms.
  \choice        Rational behaviour requires perfect foresight about the
                 future.
\end{choices}
\begin{solution}
  The assumption does not require that every actor is perfectly rational at
  every instant. Rather, in a competitive market, actors who systematically
  fail to optimise are gradually driven out (or learn to decide better),
  so that the surviving pattern of behaviour approximates rational,
  self-interested optimisation. The state, by contrast, faces different
  (often weaker) competitive pressure.
\end{solution}
\end{question}

\begin{question}[type=mc]{Simplifying Assumptions of the Basic Model}
The basic microeconomic model deliberately abstracts from real-world
frictions. Which of the following is \emph{not} one of these standard
simplifying assumptions?
\begin{choices}
  \choice        No transaction costs.
  \choice        No barriers to entry.
  \choice        Perfect information.
  \correctchoice A fixed, government-set price for every good.
\end{choices}
\begin{solution}
  The textbook model assumes away frictions such as transaction costs,
  barriers to entry, information asymmetries and time-lags — much like
  physics abstracts from air resistance when modelling free fall. A
  government-set price is not a simplifying assumption of the competitive
  model; it is itself a market intervention studied separately (price
  ceilings/floors).
\end{solution}
\end{question}


%% ═════════════════════════════════════════════════════════════════════════════
\section{Chapter 3 — Demand and Supply}
%% ═════════════════════════════════════════════════════════════════════════════

\begin{question}[type=mc]{Law of One Price}
The Law of One Price states that identical, mass-produced goods should trade
at a single price everywhere. If persistent price differences are observed
for the same good, this indicates:
\begin{choices}
  \choice        That demand for the good must be perfectly inelastic.
  \correctchoice Some form of market failure, such as barriers to arbitrage
                 or limited market access.
  \choice        That the good is a Giffen good.
  \choice        That the good is non-excludable.
\end{choices}
\begin{solution}
  If a price gap existed for an identical good, arbitrageurs would buy where
  it is cheap and sell where it is dear until the gap disappears. Persistent
  price differences therefore signal that \textbf{markets are not perfect} —
  e.g.\ due to transport costs, trade barriers, or limited information —
  which is a form of \textbf{market failure}.
\end{solution}
\end{question}

\begin{question}[type=mc]{Hotelling's Law}
Two ice-cream vendors on a beach choose where to set up. Hotelling's Law
predicts that competition drives them to:
\begin{choices}
  \choice        Locate as far apart as possible to serve different customer
                 segments.
  \correctchoice Move toward the centre of the beach and become similar to
                 one another, both in location and product.
  \choice        Randomise their location every day.
  \choice        Merge into a single monopolist.
\end{choices}
\begin{solution}
  In \textbf{Hotelling's Law}, moving toward the centre always captures more
  of the market for either vendor, so both converge to the middle of the
  beach. The same logic explains why competitors often converge toward
  similar products, similar prices, and similar political positions —
  minimal differentiation rather than maximal differentiation.
\end{solution}
\end{question}

\begin{question}[type=mc]{DEVEX, CAPEX, and OPEX}
In business-administration terminology for large investment projects, the
cost of constructing and commissioning a new asset (e.g.\ a power plant) is
called:
\begin{choices}
  \choice        DEVEX.
  \correctchoice CAPEX.
  \choice        OPEX.
  \choice        Producer surplus.
\end{choices}
\begin{solution}
  \textbf{DEVEX} (development expenditure) occurs before the final investment
  decision. \textbf{CAPEX} (capital expenditure) covers construction and
  commissioning of the asset, plus the cost of capital (interest payments)
  during repayment. \textbf{OPEX} (operational expenditure) covers ongoing
  variable costs and maintenance once the asset is running.
\end{solution}
\end{question}

\begin{question}[type=mc]{Market Forms: Buyers and Sellers}
A market with \emph{one} seller but \emph{many} buyers is called:
\begin{choices}
  \choice        An oligopoly.
  \correctchoice A monopoly.
  \choice        A monopsony.
  \choice        An oligopsony.
\end{choices}
\begin{solution}
  The market-forms grid is defined by the number of sellers (columns: one,
  few, many) and buyers (rows: one, few, many). One seller and many buyers
  is a \textbf{monopoly}. One buyer and many sellers would instead be a
  \textbf{monopsony}; few sellers give an \textbf{oligopoly}, few buyers an
  \textbf{oligopsony}.
\end{solution}
\end{question}

\begin{question}[type=mc]{Shape of the Production Possibility Frontier}
The typical Production Possibility Frontier (PPF), showing the trade-off
between producing good A and good B, is bowed outward (concave to the
origin) because:
\begin{choices}
  \choice        Factors of production are perfectly substitutable between
                 the two goods.
  \correctchoice The opportunity cost of producing more of one good rises as
                 more resources are shifted into producing it (factors are
                 not equally well suited to both uses).
  \choice        Both goods are produced with constant returns to scale and
                 identical technology.
  \choice        The firm faces a binding price ceiling.
\end{choices}
\begin{solution}
  As a producer shifts resources from good B to good A, it first reallocates
  the factors best suited to A. Further units of A require reallocating
  factors that are progressively less suited to producing A (and better
  suited to B), so the opportunity cost of each additional unit of A —
  measured in units of B forgone — \textbf{increases}. This is what gives
  the PPF its characteristic bowed-out (concave) shape.
\end{solution}
\end{question}


%% ═════════════════════════════════════════════════════════════════════════════
\section{Chapter 4 — Equilibrium and Welfare}
%% ═════════════════════════════════════════════════════════════════════════════

\begin{question}[type=mc]{Prohibitive Price and Saturation Quantity}
On a standard demand-supply diagram, the ``prohibitive price'' and the
``saturation quantity'' refer respectively to:
\begin{choices}
  \correctchoice The price so high that quantity demanded falls to zero, and
                 the quantity that would be demanded even at a price of zero.
  \choice        The equilibrium price and the equilibrium quantity.
  \choice        The price a monopolist charges, and the quantity it produces.
  \choice        The government-imposed price ceiling, and the resulting
                 excess demand.
  
\end{choices}
\begin{solution}
  The \textbf{prohibitive price} is the price at which the demand curve hits
  the price axis — nobody is willing to buy at or above it (aside from a few
  outliers). The \textbf{saturation quantity} is the quantity at which the
  demand curve hits the quantity axis at a price of zero — the amount buyers
  would take even if the good were free. Both are properties of the demand
  curve itself, distinct from the market equilibrium.
\end{solution}
\end{question}


%% ═════════════════════════════════════════════════════════════════════════════
\section{Chapter 5 — Market Disruptions}
%% ═════════════════════════════════════════════════════════════════════════════

\begin{question}[type=mc]{Shift vs.\ Pivot of the Supply Curve}
A war disrupts a specific raw material, adding a fixed cost per unit to
production, while an ad valorem import tariff is levied as a percentage of
value. On the supply-demand diagram, these two shocks:
\begin{choices}
  \choice        Both rotate the supply curve around the same point.
  \correctchoice The unit cost shock shifts the supply curve up in parallel,
                 while the ad valorem tariff pivots (rotates) it; but for
                 suitably chosen parameters both can produce the same new
                 equilibrium.
  \choice        Only affect the demand curve, not the supply curve.
  \choice        Cannot be compared, since one is a shock and the other is a
                 policy.
\end{choices}
\begin{solution}
  A per-unit cost increase shifts supply \textbf{up in parallel} (every unit
  costs the same amount more). An ad valorem tax/tariff is a percentage of
  value, so it \textbf{pivots} the supply curve (the mark-up is larger at
  higher prices). For the right combination of parameters, both can land on
  the identical new equilibrium price and quantity — meaning all the
  associated welfare (CS/PS/DWL) results carry over regardless of which
  instrument was used.
\end{solution}
\end{question}


%% ═════════════════════════════════════════════════════════════════════════════
\section{Chapter 6 — Monopolies: The Cournot Model}
%% ═════════════════════════════════════════════════════════════════════════════

\begin{question}[type=mc]{Natural Monopoly}
A ``natural monopoly'' (e.g.\ a railway network or an energy grid operator)
arises when:
\begin{choices}
  \choice        The government grants an exclusive licence for prestige
                 reasons.
  \correctchoice The fixed costs of the infrastructure are so large that a
                 single provider can supply the whole market more cheaply
                 than two or more competing providers.
  \choice        Consumers have a strong brand preference for a single firm.
  \choice        Patent protection prevents any competitor from entering.
\end{choices}
\begin{solution}
  Natural monopolies occur where \textbf{very large fixed (infrastructure)
  costs} make it inefficient to duplicate the network: a single provider's
  average cost keeps falling over the relevant range of output. Building a
  second parallel railway or grid would waste resources. This differs from
  a monopoly created by patents, brand loyalty, or state licensing.
\end{solution}
\end{question}

\begin{question}[type=mc]{Antitrust Regimes: EU vs.\ US}
Regarding the treatment of monopolies in competition law, the script notes an
important contrast:
\begin{choices}
  \choice        Both the EU and the US make simply \emph{having} a monopoly
                 illegal.
  \correctchoice In Europe, \emph{setting up} (creating) a monopoly is
                 illegal, whereas in the US, merely \emph{having} a monopoly
                 position is what antitrust law targets.
  \choice        Monopolies are legal everywhere as long as they are
                 domestic, not international.
  \choice        The EU only regulates monopolies in the energy sector.
\end{choices}
\begin{solution}
  The script highlights a legal-cultural contrast: in Europe, the emphasis is
  on preventing firms from \textbf{creating} a dominant/monopoly position in
  the first place (e.g.\ merger control), while in the US, antitrust
  enforcement can target the mere \textbf{holding} of monopoly power itself,
  regardless of how it was acquired.
\end{solution}
\end{question}


%% ═════════════════════════════════════════════════════════════════════════════
\section{Chapter 7 — Case Study: Fuel Crisis and Tax Rebate}
%% ═════════════════════════════════════════════════════════════════════════════

\begin{question}[type=mc]{Fuel Tax Rebate: Why the Price Drop Was Smaller Than
  Expected}
In spring 2026, Germany cut the fuel tax by 17 cents per litre, and Canada
cut its fuel excise tax by 10 cents per litre, yet pump prices in both
countries fell by \emph{less} than the tax cut. Given inelastic demand and
elastic supply for fuel, what would standard tax-incidence theory have
predicted?
\begin{choices}
  \choice        Prices should not have changed at all.
  \correctchoice Since demand is inelastic and supply elastic, most of the
                 tax cut should have been passed through to consumers as a
                 lower price — nearly the full amount of the rebate.
  \choice        Prices should have risen instead of falling.
  \choice        The tax rebate should have been captured entirely by the
                 government as extra revenue.
\end{choices}
\begin{solution}
  Standard incidence theory says the \textbf{less elastic} side of the market
  bears (or benefits from) more of a tax change. With demand highly
  inelastic and supply elastic, theory predicted the price should have
  fallen by almost the full amount of the rebate. In practice it fell by
  noticeably less — a sign that other frictions (shipping delays, distorted
  market structure, incomplete pass-through by retailers) were at work.
\end{solution}
\end{question}

\begin{question}[type=mc]{Restricted Supply and the Cournot Point}
The script argues that a world-wide forced reduction in available crude oil
(as in 2026) tends to push the market outcome:
\begin{choices}
  \choice        Toward the competitive equilibrium, since fewer producers
                 remain in the market.
  \correctchoice Away from the competitive equilibrium and toward the Cournot
                 point, where producer surplus is maximised.
  \choice        Toward a price of zero, since supply becomes scarce.
  \choice        Toward the saturation quantity.
\end{choices}
\begin{solution}
  When available supply is artificially restricted, remaining producers gain
  pricing power. The script notes that the resulting quantity/price
  combination can move away from the fully competitive equilibrium toward
  the \textbf{Cournot point} (the quantity that maximises producer surplus,
  as under monopoly) — which is also why some governments consider taxing
  the resulting windfall profits.
\end{solution}
\end{question}

\begin{question}[type=mc]{Comparing a Tax Cut's Cost to Its Welfare Benefit}
When a government drops a unit tax on a good whose supply has become
restricted, the case study shows that:
\begin{choices}
  \choice        The reduction in government revenue is always smaller than
                 the resulting reduction in deadweight loss.
  \correctchoice The loss of government revenue is typically \emph{larger}
                 than the reduction in deadweight loss, even though both
                 consumers and producers gain something.
  \choice        Consumer surplus increases by exactly the size of the tax
                 cut.
  \choice        Producer surplus always falls when a tax is cut.
\end{choices}
\begin{solution}
  In the case study, cutting the tax reduces the price only a little
  (because pass-through is incomplete), so consumers gain only modestly,
  while producer surplus rises more. Crucially, the government simply gives
  up tax revenue on \emph{every} unit still sold ($\Delta T \cdot x_2$),
  which exceeds the shrinkage of the deadweight-loss triangle. The exact
  split between consumers and producers depends, as always, on the ratio of
  elasticities.
\end{solution}
\end{question}


%% ═════════════════════════════════════════════════════════════════════════════
\section{Chapter 8 — Market Imperfections}
%% ═════════════════════════════════════════════════════════════════════════════

\begin{question}[type=mc]{What Can Be Patented}
Which of the following can typically \emph{not} be protected by a patent?
\begin{choices}
  \choice        A concrete technical solution to an engineering problem.
  \correctchoice A purely mathematical insight or an abstract idea.
  \choice        A newly invented industrial process.
  \choice        A specific technical device.
\end{choices}
\begin{solution}
  Patents protect concrete \textbf{technical solutions} to technical
  problems. \textbf{Ideas, mathematical insights, and abstract rules} cannot
  be patented as such. Note also the EU/Common-Law contrast between
  copyright (the sellable right to make copies) and the French/German
  \emph{droit d'auteur}/\emph{Urheberrecht} (an unsellable personal right).
\end{solution}
\end{question}

\begin{question}[type=mc]{Risk vs.\ Uncertainty}
In economic risk analysis, the key distinction between ``risk'' and
``uncertainty'' is that:
\begin{choices}
  \correctchoice Risk refers to situations where probabilities of outcomes
                 can be estimated; uncertainty refers to situations where
                 they cannot.
  \choice        Risk applies only to financial markets, uncertainty only to
                 insurance markets.
  \choice        Risk is always smaller in magnitude than uncertainty.
  \choice        There is no meaningful distinction; the terms are
                 interchangeable.
  
\end{choices}
\begin{solution}
  \textbf{Risk} describes situations that can be sorted and quantified by
  likelihood and impact — probabilities can, in principle, be estimated.
  \textbf{Uncertainty} (in the stricter, Knightian sense) refers to
  situations where outcomes are unknown and cannot be assigned meaningful
  probabilities at all.
\end{solution}
\end{question}

\begin{question}[type=mc]{The Coase Theorem in Practice}
The Coase Theorem is often cited as an argument against government
intervention in externalities. The script cautions that this is often
misapplied because, in practice:
\begin{choices}
  \choice        Property rights are always perfectly enforced.
  \correctchoice Transaction costs are never truly zero, and bargaining power
                 between the parties is often unequal.
  \choice        Externalities do not exist in real markets.
  \choice        The theorem only applies to public goods, not to
                 externalities.
\end{choices}
\begin{solution}
  The Coase Theorem's efficiency result relies on the assumption of
  \textbf{zero transaction costs} and well-defined property rights. In
  reality, negotiating, monitoring, and enforcing agreements is costly, and
  parties often have unequal bargaining power — which is why Elinor Ostrom's
  empirical work on real-world commons governance remains important, rather
  than relying on Coasean bargaining alone.
\end{solution}
\end{question}

\begin{question}[type=mc]{Procurement and the Winner's Curse}
In competitive tenders and auctions, if bidders independently estimate an
uncertain value (e.g.\ the true cost of a construction project) and the
contract goes to the lowest bidder, economic theory warns of:
\begin{choices}
  \choice        Adverse selection, because bidders raise their prices over
                 time.
  \correctchoice The winner's curse: the winning bid is likely to be based on
                 an overly optimistic (unrealistically low) estimate.
  \choice        The tragedy of the commons, because too many firms bid.
  \choice        A Pigouvian externality imposed on losing bidders.
\end{choices}
\begin{solution}
  When bidders have independent, imperfect estimates of a common, uncertain
  value, the bidder who wins is often the one who was \textbf{most
  optimistic} (i.e.\ underestimated true costs) — the \textbf{winner's
  curse}. This is why the lowest bid in a tender is not always realistic,
  and why procurement processes must account for this bias.
\end{solution}
\end{question}


%% ═════════════════════════════════════════════════════════════════════════════
\section{Chapter 9 — The State}
%% ═════════════════════════════════════════════════════════════════════════════

\begin{question}[type=mc]{Branches and Levels of Government}
A typical modern state has three institutional branches. These are:
\begin{choices}
  \choice        Federal, state (L\"ander), and communal.
  \correctchoice Executive, legislative, and judicative.
  \choice        Central bank, treasury, and parliament.
  \choice        Public law, private law, and commercial law.
\end{choices}
\begin{solution}
  The three classic branches of the state are the \textbf{executive}
  (government/administration), the \textbf{legislative} (parliament), and
  the \textbf{judicative} (courts). Separately, there are multiple
  \emph{levels} of government (federal, Länder, communal, supranational/EU)
  and a distinct central bank, which is not one of the three branches.
\end{solution}
\end{question}

\begin{question}[type=mc]{The Shell Model of Government Finance}
The ``shell model'', developed within the European System of Accounts (ESA),
is used to:
\begin{choices}
  \choice        Calculate the optimal tax rate on the Laffer curve.
  \correctchoice Make budgets of different levels of government (federal,
                 regional, social security, etc.) comparable, even when
                 governments shift items into or out of extra-budgetary
                 accounts.
  \choice        Determine which goods are public vs.\ private goods.
  \choice        Measure the Herfindahl-Hirschman Index of market
                 concentration.
\end{choices}
\begin{solution}
  The \textbf{shell model} nests the different levels and institutions of
  government (federal, state/Länder, communal, social security, government
  corporations) so that total public-sector activity can be compared across
  countries, even though governments sometimes move spending into or out of
  extra (``off-budget'') accounts for cosmetic or practical reasons.
\end{solution}
\end{question}

\begin{question}[type=mc]{Government Expenditure Ratio}
The government spending ratio (``state quota''), $SQ = G/Y$, measures:
\begin{choices}
  \choice        The share of government revenue raised through tariffs.
  \correctchoice Total government spending as a percentage of GDP.
  \choice        The ratio of federal to regional tax revenue.
  \choice        The elasticity of tax revenue with respect to the tax rate.
\end{choices}
\begin{solution}
  $SQ = G/Y$ expresses \textbf{government spending $G$ as a share of GDP
  $Y$}. It is a widely used (if imperfect) indicator of how large a role the
  state plays in an economy, and it varies substantially across countries
  with different income levels and social models.
\end{solution}
\end{question}


%% ═════════════════════════════════════════════════════════════════════════════
\section{Chapter 10 — Game Theory and Duopolies}
%% ═════════════════════════════════════════════════════════════════════════════

\begin{question}[type=mc]{Stag Hunt Game}
In the Stag Hunt game, two hunters do best if both commit to hunting stag
together, but each risks being left worse off if they commit to stag while
the other hunts hare alone. This game illustrates:
\begin{choices}
  \choice        That defection is always a dominant strategy, as in the
                 Prisoner's Dilemma.
  \correctchoice That cooperation can be the best outcome for both, but
                 requires credible commitment or trust, since both ``both
                 hunt stag'' and ``both hunt hare'' are Nash equilibria.
  \choice        That there is no Nash equilibrium in pure strategies.
  \choice        That one player always has an incentive to unilaterally
                 deviate from mutual cooperation.
  
\end{choices}
\begin{solution}
  Unlike the Prisoner's Dilemma, in the \textbf{Stag Hunt} mutual cooperation
  (both hunt stag) \emph{is} a Nash equilibrium — nobody wants to deviate
  once the other is committed. But ``both hunt hare'' is \emph{also} a Nash
  equilibrium, and it is individually safer if a hunter cannot be sure the
  other will show up. The game (going back to Rousseau) illustrates the
  importance of trust and credible commitment in achieving the
  better-for-both outcome.
\end{solution}
\end{question}

\begin{question}[type=mc]{Chicken Game}
The ``Chicken Game'' (two drivers heading toward each other, each deciding
whether to swerve) differs from the Prisoner's Dilemma mainly because:
\begin{choices}
  \choice        There is only one player in the game.
  \correctchoice The worst outcome occurs when \emph{both} players choose the
                 aggressive strategy (not swerving), rather than when both
                 choose the cooperative one.
  \choice        Both players always prefer to swerve simultaneously.
  \choice        It has no Nash equilibrium at all.
\end{choices}
\begin{solution}
  In \textbf{Chicken}, mutual ``defection'' (neither swerves) produces the
  worst possible outcome (a crash) for both, unlike the Prisoner's Dilemma
  where mutual defection is merely worse than mutual cooperation but not
  catastrophic. Chicken has two pure-strategy Nash equilibria (one swerves,
  the other doesn't) — each player wants to do the opposite of the other,
  making coordination and reputation (or commitment devices) critical.
\end{solution}
\end{question}


%% ═════════════════════════════════════════════════════════════════════════════
\section{Chapter 11 — Money and Trade}
%% ═════════════════════════════════════════════════════════════════════════════

\begin{question}[type=mc]{From Barter to Money}
In the two-good island economy (fishermen and coconut farmers), the
introduction of sea shells as a medium of exchange is used to illustrate
that:
\begin{choices}
  \choice        Money is only needed once more than 100 traders participate
                 in a market.
  \correctchoice As the variety of goods traded expands beyond a simple
                 barter pair, a common denominator (money) becomes desirable
                 to facilitate exchange.
  \choice        Barter is always more efficient than using money.
  \choice        The value of money depends only on the trust of the central
                 bank.
\end{choices}
\begin{solution}
  With only two goods (fish and coconut), direct barter can work. But as
  \textbf{commerce expands beyond a simple pair of goods}, a common medium
  of exchange (here: shells) becomes desirable so that every good does not
  need a separate exchange rate against every other good. The value of the
  medium of exchange then depends on its quantity in circulation and its
  velocity of circulation — the core idea behind the quantitative theory of
  money.
\end{solution}
\end{question}

\begin{question}[type=mc]{The Four Functions of Money}
Which of the following is \emph{not} one of the four classic functions of
money?
\begin{choices}
  \choice        Medium of exchange.
  \choice        Unit of account.
  \choice        Store of value.
  \correctchoice Guarantor of full employment.
\end{choices}
\begin{solution}
  The four classic functions of money are: \textbf{medium of exchange},
  \textbf{unit of account}, \textbf{store of value}, and \textbf{standard of
  deferred payment}. Guaranteeing full employment is a macroeconomic policy
  goal, not a function of money itself.
\end{solution}
\end{question}

\begin{question}[type=mc]{Why Do Countries Trade at All?}
Beyond comparative advantage in production costs, the script lists several
basic reasons why trade happens. Which of the following is \emph{not} one of
them?
\begin{choices}
  \choice        Acquiring goods a country otherwise would not have (due to
                 climate, resources, or know-how).
  \choice        Disposing of goods a country has in surplus (e.g.\ excess
                 capacity, dumping).
  \correctchoice Ensuring that every country's trade balance is exactly zero.
  \choice        Selling goods in order to buy other goods (the profit
                 motive).
\end{choices}
\begin{solution}
  Trade arises because countries want goods they otherwise lack (different
  climate, resources, know-how), want to get rid of surpluses, and want to
  sell in order to buy other goods (profit/maximisation motive) — echoing
  Adam Smith's argument that the wealth of nations rests on the division of
  labour. A trade balance of exactly zero is not a stated reason for trade
  and is not required for trade to be beneficial.
\end{solution}
\end{question}


%% ═════════════════════════════════════════════════════════════════════════════
\section{Chapter 12 — Protectionism}
%% ═════════════════════════════════════════════════════════════════════════════

\begin{question}[type=mc]{Non-Tariff Barriers to Trade}
Which of the following is an example of a \emph{non-tariff} barrier to
trade?
\begin{choices}
  \choice        A 10\,\% import duty on foreign cars.
  \choice        An ad valorem tax on imported electronics.
  \correctchoice A technical safety standard that foreign products struggle
                 to meet, or a lengthy customs inspection delay.
  \choice        An export subsidy paid to domestic exporters.
\end{choices}
\begin{solution}
  \textbf{Non-tariff barriers} restrict trade without directly taxing it:
  examples include technical standards, licensing and packaging/labelling
  requirements, sanitation rules, customs delays, and public-procurement
  rules that favour domestic suppliers. Tariffs, ad valorem taxes, and
  export subsidies are all \emph{price-based} instruments, not non-tariff
  barriers.
\end{solution}
\end{question}

\end{document}